%==============================================================================%
% Chapter 9 – Jasmin Implementation
%==============================================================================%

\chapter{Jasmin Implementation}
\label{ch:jasmin_impl}

\section{Overview}

The high-assurance implementation of the \haetae{} NTT is written in the
\jasmin{} programming language.  Four source/support files collectively
implement the forward NTT, inverse NTT, and base multiplication:

\begin{itemize}
  \item \texttt{reduce.jinc}: Montgomery reduction and field multiplication.
  \item \texttt{zetas.jinc}: Precomputed twiddle factor table.
  \item \texttt{poly.jinc}: Forward and inverse NTT kernels.
  \item \texttt{hpoly.jazz}: Exported entry points callable from C.
\end{itemize}

The \jasmin{} source is compiled to x86-64 assembly (\texttt{hpoly.s}) and
simultaneously extracted to an \easycrypt{} module for formal verification.
This chapter records two implementation shapes.  The original direct-loop
program is kept in the listings because it makes the correspondence with the
C reference and the \easycrypt{} loop invariants transparent.  The current
checked implementation used to regenerate \texttt{hpoly.s}, however, is the
stage-specialised full-safe constant-time version described below: the public
NTT schedule is split into fixed stage helpers, and the exported wrappers carry
explicit regular and speculative constant-time contracts.

\section{Montgomery Reduction in Jasmin}

\subsection{Source Code}

\begin{lstlisting}[language=Jasmin,
  caption={Montgomery reduction in \texttt{reduce.jinc}},
  label={lst:mont_jasmin}]
param int JPRIME = 64513;
param int JQINV  = 940508161;

inline fn __montgomery_reduce(reg u64 a) -> reg u32 {
  reg u64 t64;
  reg u32 t, r;

  t    = (32u)(a);          /* low 32 bits of a */
  t   *= (u32)(JQINV);      /* t = (a mod 2^32) * QINV mod 2^32 */
  t64  = (64s)(t);          /* sign-extend to 64 bits */
  t64 *= (u64)(JPRIME);     /* t64 = t * q (64-bit) */
  t64  = a - t64;           /* t64 = a - t*q */
  r    = (32u)(t64 >> 32);  /* r = (a - t*q) / 2^32 (arithmetic shift) */
  return r;
}
\end{lstlisting}

\subsection{Jasmin Semantics}

The \jasmin{} code uses explicit types to control the width and signedness
of arithmetic:
\begin{itemize}
  \item \jasm{(32u)(a)}: truncate 64-bit value to 32 bits (unsigned).
  \item \jasm{(u32)(JQINV)}: interpret the literal as a 32-bit unsigned value.
  \item \jasm{(64s)(t)}: sign-extend a 32-bit value to 64 bits.
  \item \jasm{(32u)(t64 >> 32)}: arithmetic right shift and truncate.
\end{itemize}

The \jasm{>>} operator on a signed 64-bit value in \jasmin{} performs an
\emph{arithmetic} right shift, matching the C behaviour for signed integers.

\subsection{Correspondence to C}

The \jasmin{} code directly implements the C reference in \Cref{lst:mont_c}:
\begin{itemize}
  \item \jasm{(32u)(a) * (u32)(JQINV)} corresponds to
        \texttt{(int32\_t)a * (int32\_t)QINV}.
  \item The sign extension \jasm{(64s)(t)} followed by multiplication by
        \jasm{JPRIME} corresponds to \texttt{(int64\_t)t * Q}.
  \item The arithmetic right shift \jasm{>> 32} corresponds to
        \texttt{>> 32} on a signed value.
\end{itemize}

\section{Field Multiplication}

\begin{lstlisting}[language=Jasmin,
  caption={Field multiplication in \texttt{reduce.jinc}},
  label={lst:fqmul_jasmin}]
inline fn __fqmul(reg u32 a, reg u32 b) -> reg u32 {
  reg u64 t64;
  reg u32 r;

  t64 = (64s)(a);
  t64 *= (64s)(b);
  r   = __montgomery_reduce(t64);
  return r;
}
\end{lstlisting}

This computes $a \cdot b \cdot R^{-1} \bmod q$.  The refinement proof
discharges the product-range side condition for the NTT cases, including
16-bit twiddle factors multiplied by coefficients bounded by the current
stage invariant.

\section{Forward NTT Kernel}

\subsection{Original Direct Source Code}

\begin{lstlisting}[language=Jasmin,
  caption={Original direct forward NTT kernel in \texttt{poly.jinc}},
  label={lst:ntt_jasmin}]
inline fn _poly_ntt(reg mut ptr u32[256] rp) -> reg ptr u32[256] {
  reg u32 zeta, t;
  reg int len, start, j, zetasctr;

  zetasctr = 0;
  len = 128;

  while (len > 0) {
    start = 0;
    while (start < 256) {
      zetasctr += 1;
      zeta = (u32) jzetas[zetasctr];

      j = start;
      while (j < start + len) {
        t = __fqmul(zeta, rp[j + len]);

        rp[j + len] = rp[j] - t;
        rp[j]       = rp[j] + t;
        j += 1;
      }
      start = j + len;
    }
    len >>= 1;
  }
  return rp;
}
\end{lstlisting}

\Cref{lst:ntt_jasmin} is the original direct implementation.  It is retained
as the readable algorithmic baseline: it exposes the same three nested loops as
the C reference.  In the current compiler-checked source, this public outer
schedule is no longer represented by run-time variables \jasm{len},
\jasm{start}, and \jasm{zetasctr}.  Instead, each stage is a separate helper
with literal bounds and a literal twiddle-table base.

\subsection{Structural Correspondence}

The \jasmin{} kernel \texttt{\_poly\_ntt} has a direct structural
correspondence to the C reference (\Cref{lst:c_ntt}) and the \easycrypt{}
specification (\Cref{lst:ntt_spec}):

\begin{center}
\small
\begin{tabular}{
  >{\raggedright\arraybackslash}p{0.18\textwidth}
  >{\raggedright\arraybackslash}p{0.42\textwidth}
  >{\raggedright\arraybackslash}p{0.28\textwidth}}
\toprule
Aspect & C reference & \jasmin{} \\
\midrule
Outer loop & \texttt{len = 128; len > 0; len >>= 1} &
             \texttt{while (len > 0)} \\
Middle loop & \texttt{start < 256; start += 2*len} &
              equivalent \texttt{while} \\
Inner loop & \texttt{j = start; j < start+len; j++} &
             equivalent \texttt{while} \\
Twiddle & \texttt{zetas[k++]} & \texttt{zetasctr += 1; jzetas[zetasctr]} \\
Butterfly & \texttt{mont\_reduce(zeta * a[j+len])} & same \\
\bottomrule
\end{tabular}
\normalsize
\end{center}

The key difference is the use of explicit type casts in \jasmin{} to manage
word widths.

\subsection{Array Access Pattern}

In \jasmin{}, the global constant table \texttt{jzetas} is accessed as a
\jasm{u32} array.  The entry at index $k$ is the Montgomery form of the
$k$-th twiddle factor.  The access pattern increments \jasm{zetasctr}
before each read, so the forward kernel consumes \jasm{jzetas[1]} through
\jasm{jzetas[255]}; \jasm{jzetas[0]} is padding.

\subsection{Full-Safe Stage-Specialised Forward Kernel}

The full-safe implementation keeps the same butterfly operations and twiddle
order, but makes every stage bound explicit to the \jasmin{} safety checker.
For the first forward stage, the checker sees the literal facts
\jasm{block < 1}, \jasm{j < 128}, and \jasm{offset = idx + 128}; hence both
\jasm{rp[idx]} and \jasm{rp[offset]} are statically within the 256-word
polynomial.  The remaining helpers instantiate the same pattern for
\(64,32,16,8,4,2,1\).

\begin{lstlisting}[language=Jasmin,
  caption={Current full-safe forward NTT stage shape in \texttt{poly.jinc}},
  label={lst:ntt_jasmin_safe_stage}]
fn _poly_ntt_stage_128(reg mut ptr u32[256] rp)
    -> reg ptr u32[256] {
  reg ptr u32[256] zetasp;
  reg ui64 block, j, start, idx, offset;
  reg u32 zeta, t, s, coeff;

  zetasp = jzetas;
  block = 0;
  while (block < 1) {
    zeta = zetasp[1 + block];
    start = block;
    start <<= 8;
    j = 0;
    while (j < 128) {
      idx = j;
      idx += start;
      s = rp[idx];
      offset = idx;
      offset += 128;
      coeff = rp[offset];
      t = __fqmul(zeta, coeff);
      coeff = s;
      coeff -= t;
      rp[offset] = coeff;
      s += t;
      rp[idx] = s;
      j += 1;
    }
    block += 1;
  }
  return rp;
}

fn _poly_ntt(reg mut ptr u32[256] rp) -> reg ptr u32[256] {
  rp = _poly_ntt_stage_128(rp);
  rp = _poly_ntt_stage_64(rp);
  rp = _poly_ntt_stage_32(rp);
  rp = _poly_ntt_stage_16(rp);
  rp = _poly_ntt_stage_8(rp);
  rp = _poly_ntt_stage_4(rp);
  rp = _poly_ntt_stage_2(rp);
  rp = _poly_ntt_stage_1(rp);
  return rp;
}
\end{lstlisting}

The transformation is proof-preserving at the algorithmic level: the sequence
of helper calls is the fixed unrolling of the original public \jasm{len}
schedule.  The implementation change is motivated by compiler verification,
not by a change in the mathematical NTT.

\section{Inverse NTT Kernel}

\begin{lstlisting}[language=Jasmin,
  caption={Original direct inverse NTT kernel in \texttt{poly.jinc} (condensed)},
  label={lst:invntt_jasmin}]
inline fn _poly_invntt(reg mut ptr u32[256] rp) -> reg ptr u32[256] {
  reg u32 zeta, t, coeff;
  reg int len, start, j, zetasctr;

  zetasctr = 0;
  len = 1;

  while (len < 256) {
    start = 0;
    while (start < 256) {
      zeta = (u32) jzetas_inv[zetasctr];
      zetasctr += 1;

      j = start;
      while (j < start + len) {
        t           = rp[j];
        coeff       = rp[j + len];
        rp[j]       = t + coeff;
        rp[j + len] = __fqmul(zeta, t - coeff);

        j += 1;
      }
      start = j + len;
    }
    len <<= 1;
  }

  /* Final scaling */
  zeta = (u32) jzetas_inv[255];
  j = 0;
  while (j < 256) {
    rp[j] = __fqmul(zeta, rp[j]);
    j += 1;
  }
  return rp;
}
\end{lstlisting}

The inverse butterfly loop consumes \texttt{jzetas\_inv[0..254]}.  The final
scaling loop is the only place where \texttt{jzetas\_inv[255]} is read.
\Cref{lst:invntt_jasmin} is likewise retained as the original direct-loop
form.  The current checked source splits the inverse transform into helpers
\jasm{\_poly\_invntt\_stage\_1}, \jasm{\_poly\_invntt\_stage\_2}, \(\ldots\),
\jasm{\_poly\_invntt\_stage\_128}, followed by \jasm{\_poly\_invntt\_scale}.
The helper sequence fixes the twiddle bases
\(0,128,192,224,240,248,252,254\), and the scale helper reads
\jasm{jzetas\_inv[255]}.

\subsection{Inverse Butterfly Structure}

The inverse butterfly in \jasmin{} computes:
\begin{align*}
  t &\gets \texttt{rp}[j] \\
  \texttt{rp}[j] &\gets t + \texttt{rp}[j+\ell] \\
  \texttt{rp}[j+\ell] &\gets \op{mont}(\zeta \cdot (t - \texttt{rp}[j+\ell]))
\end{align*}

This is the inverse (Gentleman--Sande) butterfly.  It is the algebraic inverse
of the Cooley--Tukey butterfly: if the forward butterfly computes $(a+b, a-b)$
with twiddle $w$ applied to the lower output, the inverse butterfly computes
$(a+b, w(a-b))$ and is composed with the twiddle application on the upper
output in the forward case.

\begin{lstlisting}[language=Jasmin,
  caption={Current full-safe inverse NTT schedule in \texttt{poly.jinc}},
  label={lst:invntt_jasmin_safe_schedule}]
fn _poly_invntt(reg mut ptr u32[256] rp) -> reg ptr u32[256] {
  rp = _poly_invntt_stage_1(rp);
  rp = _poly_invntt_stage_2(rp);
  rp = _poly_invntt_stage_4(rp);
  rp = _poly_invntt_stage_8(rp);
  rp = _poly_invntt_stage_16(rp);
  rp = _poly_invntt_stage_32(rp);
  rp = _poly_invntt_stage_64(rp);
  rp = _poly_invntt_stage_128(rp);
  rp = _poly_invntt_scale(rp);
  return rp;
}
\end{lstlisting}

\section{Exported Entry Points}

\begin{lstlisting}[language=Jasmin,
  caption={Original exported functions in \texttt{hpoly.jazz}},
  label={lst:hpoly}]
export fn poly_ntt_jazz(reg mut ptr u32[256] rp)
    -> reg ptr u32[256] {
  rp = _poly_ntt(rp);
  return rp;
}

export fn poly_invntt_jazz(reg mut ptr u32[256] rp)
    -> reg ptr u32[256] {
  rp = _poly_invntt(rp);
  return rp;
}

export fn poly_basemul_jazz(
    reg mut ptr u32[256] rp,
    reg ptr u32[256] ap,
    reg ptr u32[256] bp)
    -> reg ptr u32[256] {
  rp = _poly_basemul(rp, ap, bp);
  return rp;
}
\end{lstlisting}

\Cref{lst:hpoly} is the original wrapper form.  The current full-safe
constant-time wrapper form adds explicit contracts for both ordinary
constant-time checking and speculative constant-time checking.  The
\jasm{\#init\_msf()} call initializes the microarchitectural speculation
fence state required by \texttt{jasmin-ct --speculative}; in the generated
assembly it emits an \texttt{lfence} at the public entry point.

\begin{lstlisting}[language=Jasmin,
  caption={Current constant-time exported wrappers in \texttt{hpoly.jazz}},
  label={lst:hpoly_safe}]
#[ct = "rp -> rp"]
#[sct = "{ ptr: transient, val: secret } -> { ptr: public, val: secret }"]
export
fn poly_ntt_jazz(reg mut ptr u32[256] rp) -> reg ptr u32[256] {
  _ = #init_msf();
  rp = _poly_ntt(rp);
  return rp;
}

#[ct = "rp -> rp"]
#[sct = "{ ptr: transient, val: secret } -> { ptr: public, val: secret }"]
export
fn poly_invntt_jazz(reg mut ptr u32[256] rp) -> reg ptr u32[256] {
  _ = #init_msf();
  rp = _poly_invntt(rp);
  return rp;
}

#[ct = "rp × ap × bp -> {ap, bp, rp}"]
#[sct = "
{ ptr: transient, val: secret } ×
{ ptr: transient, val: secret } ×
{ ptr: transient, val: secret } ->
{ ptr: public, val: secret }
"]
export
fn poly_basemul_jazz(reg mut ptr u32[256] rp,
                     reg const ptr u32[256] ap bp) -> reg ptr u32[256] {
  _ = #init_msf();
  rp = _poly_basemul(rp, ap, bp);
  return rp;
}
\end{lstlisting}

The \jasm{export} keyword triggers both assembly generation and \easycrypt{}
extraction.  The ABI of \jasm{poly\_ntt\_jazz} takes a pointer to a 256-word
array in a register and returns the same pointer after in-place modification.
The argument names follow the conventional pointer suffix \jasm{p}.  Thus
\jasm{rp} is the result-polynomial pointer.  It is mutable and is used in-place
by \jasm{poly\_ntt\_jazz} and \jasm{poly\_invntt\_jazz}; in
\jasm{poly\_basemul\_jazz} it names the destination array that receives the
coefficientwise products.  The names \jasm{ap} and \jasm{bp} are the two
read-only input-polynomial pointers for the base-multiplication operands
\(a\) and \(b\).  All three pointers range over arrays of
\jasm{u32[256]} coefficients, but only \jasm{rp} is declared \jasm{mut};
the current checked wrapper declares \jasm{ap} and \jasm{bp} as
\jasm{const} to make their read-only role explicit.

\section{Safety and Constant-Time Verification}

The \jasmin{} verification obligations attached to this chapter are not only
functional.  Functional refinement, which is carried out later in
\easycrypt{}, proves that the extracted procedures implement the mathematical
NTT.  The compiler-side checks discussed here prove two lower-level
properties that must hold before the functional statement can be trusted as an
implementation statement: the program must be \emph{safe}, so that the
low-level execution is defined, and it must be \emph{constant time}, so that
its control flow and memory-access pattern do not depend on polynomial
coefficients.  The account below follows the local \jasmin{} checkout at
\texttt{\$HOME/jasmin}: the safety checker is implemented under
\texttt{compiler/safetylib}, the regular constant-time checker in
\texttt{compiler/src/ct\_checker\_forward.ml}, the speculative checker in
\texttt{compiler/src/sct\_checker\_forward.ml}, and the command-line entry
point for the latter two in \texttt{compiler/entry/jasmin\_ct.ml}.
The conceptual background is the CCS 2017 \jasmin{} paper~\cite{Almeida2017},
which presents the language as a way to keep assembly-level control while
making source-level safety and timing analyses tractable.  That paper describes
the original source-analysis pipeline in terms of two translations: a
functional embedding used to prove safety, and a constant-time instrumentation
whose product program reduces timing security to a safety proof.  The local
tooling used in this dissertation has evolved, but the proof principles are the
same: first establish that the execution is defined, then establish that the
observable trace is independent of secrets, and finally rely on the
\jasmin{} compiler's preservation argument to carry those source-level
properties to the generated assembly.

\subsection{What Safety Means for This Code}

Safety is the definedness layer of the proof.  The \jasmin{} type system gives
each pointer and array an element width and a length, for example
\jasm{reg mut ptr u32[256] rp}.  The safety checker then verifies that every
operation which could be undefined at the assembly level is executed only
under conditions that make it valid.  For this scalar NTT implementation the
important safety obligations are:
\begin{itemize}
  \item \textbf{Valid memory accesses.}  Every read and write through
        \jasm{rp}, \jasm{ap}, or \jasm{bp} must be inside the corresponding
        256-word polynomial object.  Every access to \jasm{jzetas} or
        \jasm{jzetas\_inv} must be inside the 256-entry global table.
  \item \textbf{Alignment and access width.}  A \jasm{u32} access must be made
        at an address that is valid for a 32-bit load or store.  In the source
        this is expressed by typed \jasm{u32[256]} pointers rather than by
        byte-level address arithmetic.
  \item \textbf{Defined control flow.}  Loops must have public, bounded
        schedules.  The checker reports termination obligations separately
        from memory obligations; in this implementation all loops count over
        fixed stage constants.
  \item \textbf{Defined use of primitive operations.}  Word arithmetic itself
        is not C-style undefined arithmetic: \jasmin{} word operations have
        fixed bit-vector semantics.  The safety proof therefore does not need
        to prove absence of signed overflow in Montgomery arithmetic.  It does
        need to prove that primitive operations are used at valid types and
        widths and that memory operands are valid.
\end{itemize}

In the CCS 2017 presentation, this safety condition is explained through a
near one-to-one functional embedding of \jasmin{} into a verification language:
fixed-size \jasmin{} arrays become fixed-size sequences, memory is modeled as
byte blocks, and every memory read or write is guarded by validity assertions
for the addressed byte range~\cite{Almeida2017}.  The modern
\jasm{-checksafety} command no longer needs to be described as a Dafny proof in
order to use it in this artifact, but the generated obligations have the same
shape.  A safety success means that every array index, memory address, shift
amount, and other potentially failing operation has been proved valid under the
function's typed calling convention and loop invariants.

The present implementation uses a deliberately small fragment of that memory
model.  The source code in \texttt{poly.jinc} has no explicit raw memory
access of the form \jasm{[p + e]} and no source-level stack array.  Its
observable memory operations are typed array/pointer operations:
\jasm{rp[idx]}, \jasm{rp[offset]}, \jasm{ap[i]}, \jasm{bp[i]}, and reads from
the global tables \jasm{jzetas} and \jasm{jzetas\_inv}.  Thus the abstract
byte-block obligations from CCS17 specialize here to simple word-array
obligations: each caller-provided polynomial pointer must designate a valid
256-word \jasm{u32} region, and every computed index must lie in
\([0,256)\).  This is why the implementation avoids pointer arithmetic in the
source even though the generated x86 code necessarily lowers the same facts to
base-plus-scaled-index addressing.

The mathematical safety argument for the original direct-loop code is simple.
For a forward stage with length \(\ell\), the inner loop reads and writes
\(\texttt{rp}[j]\) and \(\texttt{rp}[j+\ell]\), with
\[
  j \in [\texttt{start}, \texttt{start}+\ell)
  \quad\text{and}\quad
  j+\ell \in [\texttt{start}+\ell,\texttt{start}+2\ell).
\]
The middle loop advances \(\texttt{start}\) by \(2\ell\), so the two half-blocks
are always contained in one interval
\([\texttt{start},\texttt{start}+2\ell)\subseteq[0,256)\).  The inverse
butterfly has the same address shape.  The twiddle-table argument is also
fixed: the forward transform consumes \(\texttt{jzetas}[1]\) through
\(\texttt{jzetas}[255]\); the inverse butterfly stages consume
\(\texttt{jzetas\_inv}[0]\) through \(\texttt{jzetas\_inv}[254]\); and the
final scaling loop consumes \(\texttt{jzetas\_inv}[255]\).
Base multiplication has the simplest safety shape.  The loop counter
\jasm{i} ranges from \(0\) to \(255\), each iteration reads exactly
\jasm{ap[i]} and \jasm{bp[i]}, and it writes exactly \jasm{rp[i]}.  The
declarations in \texttt{hpoly.jazz} mark \jasm{ap} and \jasm{bp} as
\jasm{const} pointers, so the safety proof for this routine also matches the
intended data-flow role: the two source polynomials are read, and the result
polynomial is the only destination.

The stage-specialised implementation exists to make these same facts directly
visible to the automatic checker.  A human can infer from the direct program
that the triple \((\texttt{len},\texttt{start},\texttt{j})\) maintains the
range invariant above.  An abstract interpreter is deliberately more local and
more conservative: it must discover enough arithmetic relationships between
mutable loop variables to prove every array access.  The full-safe source
therefore replaces the public run-time stage variable \jasm{len} by eight
helper functions with literal lengths \(128,64,32,16,8,4,2,1\).  In a helper
with length \(\ell\), the code has the schematic form
\[
  \texttt{block}<\frac{256}{2\ell},\qquad
  \texttt{j}<\ell,\qquad
  \texttt{start}=\texttt{block}\cdot 2\ell,\qquad
  \texttt{offset}=\texttt{start}+\texttt{j}+\ell.
\]
Consequently the checker only needs local arithmetic to derive
\[
  0 \leq \texttt{idx} < 256
  \quad\text{and}\quad
  0 \leq \texttt{offset} < 256.
\]
The first forward helper shown in \Cref{lst:ntt_jasmin_safe_stage} is the
extreme case: \(\texttt{block}<1\), \(\texttt{j}<128\), and
\(\texttt{offset}=\texttt{j}+128\).  Later helpers instantiate the same proof
with smaller \(\ell\) and more blocks.  The inverse helpers use the same
address proof in the opposite stage order, and the final scaling helper has
only the single index obligation \(0\leq j<256\).

\subsection{Regular Constant-Time Typing}

The regular constant-time checker is a security type checker.  Its lattice has
public, secret, and polymorphic levels.  Public data may influence observable
timing behaviour; secret data may not.  Polymorphic levels are used in
function contracts so that a returned value can be described as depending on a
particular input without forcing that input to be public.  The checker then
walks the program and enforces the following discipline.

The semantic property behind this type discipline is trace equality.  The
\jasmin{} paper instruments the operational semantics with a leakage trace
recording the branches taken and the memory read/write operations performed
during execution.  A command is constant time when any two executions that
agree on public inputs produce the same leakage trace, even if their secret
state differs~\cite{Almeida2017}.  This is a termination-insensitive property,
so safety is checked separately and first; the dissertation follows the same
ordering by pairing \jasm{jasminc -checksafety} with \jasm{jasmin-ct}.

\begin{itemize}
  \item Conditions of \jasm{if} statements, \jasm{while} guards, and loop
        bounds must type-check as public.  Thus a secret coefficient cannot
        determine whether an instruction executes.
  \item Memory addresses and array indices must type-check as public.  Thus a
        secret coefficient cannot determine which cache line is read or
        written.
  \item Operations classified by the architecture as not constant time may be
        used only on public operands.  This is why the checker is parameterised
        by the architecture-specific predicate passed to
        \texttt{jasmin-ct}.
  \item The result of a memory load is treated conservatively as secret unless
        its source is known more precisely through the typed array discipline.
        This is exactly the desired default for cryptographic state.
\end{itemize}

This is a proof relative to the \jasmin{} x86 leakage model used by the
checker.  Following CCS17, the observable trace records branches and memory
read/write operations, while instructions classified as constant-time by the
architecture model do not add secret-dependent leakage obligations.  This
matters for the present code because Montgomery multiplication lowers to
\texttt{imul} and \texttt{sar} instructions on coefficient-derived registers:
the checker accepts these arithmetic operations under the selected x86 model,
but it would reject a use of secret data in an operation that the model marks
as non-constant-time or in any address/branch expression.

The CCS 2017 toolchain checked this property by inserting publicness
annotations on control-flow conditions, memory accesses, and array accesses,
then using a Boogie product program that executes two copies of the program in
lockstep.  Publicness assertions become relational equalities between the
original and shadow executions; proving the product program safe proves that
the two leakage traces coincide~\cite{Almeida2017}.  The local
\jasm{jasmin-ct} checker used here is a lighter direct type checker, but it is
enforcing the same relational idea syntactically: anything that contributes to
the leakage trace must be public.

For the \haetae{} NTT, all data-dependent values are coefficient values or
intermediate products inside \jasm{\_\_fqmul} and
\jasm{\_\_montgomery\_reduce}.  Those values flow through additions,
subtractions, casts, shifts, and multiplications, but they never flow into a
branch condition, a loop bound, or an address expression.  The schedule is
fixed by the public stage constants.  The table indices are also public:
\jasm{1 + block}, \jasm{2 + block}, \(\ldots\), and the inverse-stage bases are
literal constants.  Therefore the memory trace of the NTT is independent of
the coefficients even though the loaded and stored words themselves are
secret.

The same trace argument applies to \jasm{\_poly\_basemul}.  The routine has
one public loop counter \jasm{i}; on each iteration it reads \jasm{ap[i]} and
\jasm{bp[i]}, computes \jasm{\_\_fqmul(a, b)}, and writes the result to
\jasm{rp[i]}.  The loaded coefficients \jasm{a} and \jasm{b} are secret, and
the Montgomery product is secret, but the only branch condition is
\jasm{i < HAETAE\_N} and the only addresses use the public index \jasm{i}.
The generated assembly consequently contains the expected public-counter loop:
loads from the two source bases at scale four, a store to the result base at
the same public index, and a comparison of the loop counter with \(256\).  This
is exactly the CCS17 leakage discipline instantiated on the actual scalar
base-multiplication code.

The exported contracts in \Cref{lst:hpoly_safe} record this policy at the ABI
boundary.  The contract \jasm{\#[ct = "rp -> rp"]} says that the output of the
in-place NTT has the same security dependence as its input polynomial.  It
does not say that the polynomial is public; it says that the function preserves
the input's security level while the checker separately enforces public use of
control and address expressions.  For base multiplication, the contract
\jasm{\#[ct = "rp * ap * bp -> \{ap, bp, rp\}"]} records
the analogous three-argument policy.  Here the identifiers \jasm{rp},
\jasm{ap}, and \jasm{bp} are also used as named security levels in the
\jasmin{} annotation syntax.  They are chosen to match the pointer variables:
\jasm{rp} labels the result-pointer argument, while \jasm{ap} and \jasm{bp}
label the two source-polynomial arguments.  The result type
\jasm{\{ap, bp, rp\}} is intentionally conservative: the returned pointer is
allowed to carry dependence on all three input labels, even though the source
body returns the destination pointer after writing all 256 cells.  The
important constant-time fact is separate from minimality of this return label:
the checker still requires the loop index used in \jasm{ap[i]}, \jasm{bp[i]},
and \jasm{rp[i]} to be public, so secret coefficient values can affect only the
words written to \jasm{rp}, not which locations are read or written.

\subsection{Speculative Constant-Time Typing}

The speculative checker strengthens the regular property against transient
execution.  In the source language this is expressed by the \jasm{sct}
attribute.  Unlike the regular \jasm{ct} contract, an \jasm{sct} type
distinguishes the security of a pointer value from the security of the memory
contents reachable through that pointer:
\[
  \text{\jasm{\{ ptr: }}\ell_p\text{\jasm{, val: }}\ell_v
  \text{\jasm{ \}}}.
\]
It also distinguishes ordinary execution from speculative execution.  The
special level \jasm{transient} is printed by the Jasmin checker for the pair
\((\text{normal public},\text{speculative secret})\).  Thus an argument of type
\[
  \jasm{\{ ptr: transient, val: secret \}}
\]
means: on the architecturally committed path, the pointer is public enough to
be used as an address; during speculative execution, the checker treats it as
potentially secret until a speculation barrier or hardening operation has
established a public status.  The memory values remain secret in both modes.

The instruction \jasm{\#init\_msf()} is the explicit boundary used by the
exported wrappers.  In the checker it initializes the \jasmin{} MSF state and
normalises speculative security levels after the public entry point has been
reached.  In the x86-64 lowering used here, the corresponding exported
assembly contains an \texttt{lfence}.  The wrappers therefore have the shape
\[
  \text{transient pointer input}
  \xrightarrow{\ \text{\jasm{\#init\_msf()}}\ }
  \text{public pointer schedule with secret values}.
\]
After that boundary, the same fixed-stage argument used for regular
constant-time applies, but now every array access is checked in the stronger
speculative type environment.  The output contract
\jasm{\{ ptr: public, val: secret \}} states that returning the pointer does
not leak secret information and that the polynomial contents are still secret.
For \jasm{poly\_basemul\_jazz}, the three input arguments all have the
same speculative type: each pointer is transient at entry and each pointed-to
polynomial value is secret.  After \jasm{\#init\_msf()}, the checker may use
\jasm{rp}, \jasm{ap}, and \jasm{bp} as public address bases, while the values
loaded from \jasm{ap[i]} and \jasm{bp[i]} remain secret and flow only into the
stored coefficient \jasm{rp[i]}.

This distinction matters for an exported cryptographic routine.  A regular
constant-time proof rules out leaks from committed branches and committed
addresses.  Speculative constant-time additionally prevents the proof from
silently assuming that a public-looking pointer or bound remains public under
transient execution before the entry fence has taken effect.  In this
implementation there are no secret-dependent bounds to harden inside the
polynomial kernels; the only required hardening point is the public wrapper
entry.

\subsection{Why Source-Level Checks Apply to Assembly}

The \jasmin{} paper separates two questions that are easy to conflate.  The
first is whether the source program is safe and constant time.  The second is
whether compilation can invalidate those properties.  Jasmin's answer is to
restrict the compiler to predictable transformations and prove semantic
preservation in Coq, then argue that the compiler also preserves the
constant-time leakage property~\cite{Almeida2017}.  In particular, the compiler
does not perform the aggressive optimisations that are dangerous for
constant-time code, such as replacing computations on secret values with table
lookups or introducing case splits on secret-dependent ranges.  Conditional
operations are lowered to constant-time conditional instructions rather than to
secret-dependent branches, and transformations such as loop unrolling or
constant propagation remove only public, fixed control structure.

This point is central for the present artifact.  The verification commands are
run on \texttt{jasmin/hpoly.jazz}, not by manually auditing every generated
assembly instruction.  That is meaningful because the \jasmin{} compilation
discipline is designed so that source-level public schedule and public address
facts remain public in the assembly trace.  The generated \texttt{hpoly.s}
therefore inherits the fixed-stage NTT schedule: any additional loads or stores
introduced by allocation are at locations determined by the public calling
convention and public loop counters, not by coefficient values.

This can be seen directly in the generated scalar assembly.  The exported
wrappers begin with an \texttt{lfence} corresponding to \jasm{\#init\_msf()}.
The base-multiplication loop lowers \jasm{ap[i]}, \jasm{bp[i]}, and
\jasm{rp[i]} to indexed memory operands with the public loop counter as the
scaled index; the loop branch compares that counter against the constant
\(256\).  Likewise, each NTT stage loop branches on public counters and uses
addresses formed from the public base pointer, public stage constants, and
public counters.  Secret coefficients appear in arithmetic registers and in
stored words, but not in branch predicates or address registers.

\subsection{Instantiation for the \haetae{} NTT}

Combining the safety and timing arguments gives a precise invariant for every
exported routine:
\begin{itemize}
  \item \textbf{Public schedule.}  The number of iterations in every loop is a
        function only of the fixed transform size \(256\) and the current
        stage length.  It is independent of coefficient values.
  \item \textbf{Public addresses.}  The address expressions for
        \jasm{rp[idx]}, \jasm{rp[offset]}, \jasm{ap[i]}, \jasm{bp[i]}, and
        the twiddle tables are affine expressions in public loop counters and
        literal stage constants.
  \item \textbf{Secret arithmetic.}  Coefficients, products, Montgomery
        reductions, and butterfly sums are allowed to be secret.  They affect
        stored values but never affect which instructions execute or which
        addresses are touched.
  \item \textbf{Defined accesses.}  The same public affine expressions are
        bounded strongly enough to prove all array indices lie in
        \([0,256)\).
  \item \textbf{Speculation boundary.}  Each exported wrapper initializes the
        MSF state before calling the internal kernel, turning transient pointer
        inputs into public pointer schedules while preserving secret values.
\end{itemize}

The implementation style is therefore intentionally conservative.  It avoids
data-dependent conditional subtraction, table lookups indexed by coefficients,
variable-length loops, and secret-dependent pointer arithmetic.  Some of these
patterns could be made secure with additional masking or hardening, but they
would enlarge the trusted timing argument.  The scalar helperized code keeps
the checker's proof obligation identical to the informal cryptographic rule:
only the public NTT schedule determines the trace.

This is the main reason the verification claim is persuasive rather than only
a tool invocation.  The desired statement is a two-run noninterference
property: two calls that satisfy the same public ABI and pointer-validity
preconditions, but contain different secret coefficient values, must have the
same branch and memory-address trace.  Safety supplies the defined-execution
premise for those two runs; regular constant-time typing proves that every
expression contributing to the trace is public; speculative constant-time
typing proves that the exported entry boundary first converts transient pointer
inputs into public address bases without making polynomial contents public;
and the Jasmin compiler-preservation theorem carries those source facts to
the generated assembly.  In this implementation each premise can be inspected
directly in the source: loop guards are public constants and counters, address
expressions are affine in public counters and stage offsets, and coefficients
appear only in arithmetic registers and stored words.  The checker is
therefore validating the same invariant a manual audit would look for, but it
does so exhaustively over all stages, all loop iterations, and all exported
routines.

\subsection{Concrete Checks}

The following checks were run on \texttt{jasmin/hpoly.jazz} after regenerating
the source:

\begin{lstlisting}[
  caption={Compiler checks for the full-safe constant-time implementation},
  label={lst:jasmin_safe_ct_checks}]
jasminc -checksafety -o /tmp/hpoly-safe.s jasmin/hpoly.jazz
  poly_ntt_jazz:    No Safety Violation
  poly_invntt_jazz: No Safety Violation
  poly_basemul_jazz: No Safety Violation

jasmin-ct jasmin/hpoly.jazz
  regular constant-time check: passed

jasmin-ct --speculative jasmin/hpoly.jazz
  speculative constant-time check: passed
\end{lstlisting}

The three results should be read together.  The safety result says that the
compiler found no possible out-of-bounds, invalid, or misaligned execution of
the exported routines under the typed input-pointer assumptions.  The regular
constant-time result says that committed control flow and committed memory
addresses are independent of secret coefficients.  The speculative
constant-time result says that this remains valid in the stronger transient
model induced by the \jasm{sct} annotations and the \jasm{\#init\_msf()}
boundary.  Consequently, the checked assembly has a fixed instruction and
memory-access schedule for all coefficient values, while the values carried by
that schedule remain secret.

\section{Generated Assembly}

The compiled assembly \texttt{hpoly.s} contains straightforward x86-64 code.
Key assembly features include:
\begin{itemize}
  \item Register allocation: stage-local variables (\texttt{block},
        \texttt{start}, \texttt{j}, \texttt{idx}, and \texttt{offset}) are
        assigned to general-purpose registers.
  \item Montgomery reduction uses \texttt{imul} (signed 64-bit multiply)
        and \texttt{sar} (arithmetic shift right) instructions.
  \item The exported wrappers contain an entry \texttt{lfence}, emitted by
        \jasm{\#init\_msf()}, to satisfy the speculative constant-time
        contract.
  \item No SIMD instructions are used in the scalar implementation; a
        vectorised (AVX2) version is possible but beyond the scope of this
        dissertation.
\end{itemize}

The correspondence between the \jasmin{} source, the generated assembly, and
the \easycrypt{} extraction is guaranteed by the \jasmin{} compiler
correctness theorem~\cite{Almeida2017}, which is itself a machine-checked
proof in the Coq proof assistant.
